% !TeX root = ../xdupgthesis_template_lc.tex

岁月不居，时节如流。转眼间，三年硕士求学生涯已近尾声。从初入校园时的青涩与忐忑，到如今在学习、科研与生活中逐渐沉淀与成长，这段时光既有面对困难时的迷茫与压力，也有不断探索后的收获与欣慰。临近毕业，心中并不只有学成毕业的轻松，更多的是对这段求学经历以及一路同行之人的珍视与不舍。

饮水思源，师恩难忘。首先，我要衷心感谢我的导师毛国强老师。毛老师始终给予我悉心指导与耐心帮助，每当我在科研过程中遇到困难与困惑时，毛老师总能从更高的视角给予我启发，使我逐渐学会以更加严谨的态度思考和解决问题。其严谨踏实的治学态度和认真负责的工作作风，也将长期影响我今后的学习与生活。此外，亦要感谢任效江老师，任老师在求学期间给予我许多帮助，使我在困惑之时得以拨云见日。在此，再次感谢两位老师对我的帮助与指导。

同窗共勉，温暖常在。我也要感谢在科研工作中给予我帮助的同门与师兄师姐。感谢刘烊烊师兄、杨琦玥师姐、黄研师兄在求职择业过程中给予我的经验分享与建议，让我在毕业选择面前少了一些迷茫，多了一些从容。感谢柳文崇师兄、向倪竺师姐在课题学习和问题分析过程中的耐心指导，在我遇到疑惑与困难时总能及时点拨，使我少走了许多弯路。感谢毛健强在数据采集和实验推进中给予我的协助，那些在炎热天气下共同完成采集任务的经历，让我深切体会到科研工作的艰辛与严谨。感谢肖磊、邱嘉乐、翟肖童、张溢泉、梁启辉在学习与生活中的陪伴，那些一起聚餐、出行的轻松时刻，构成了研究生生活中最珍贵的记忆。同时，我也要感谢本科便认识的徐毅洁同学，感谢她在许多重要而艰难的时刻相互扶持、共渡难关。

感恩相遇，幸得同行。我要感谢 309 实验室这一温暖的集体。感谢付天烜、王科银两位博士师兄以及侯照莹在学习和生活中给予我的帮助与关照。无论是学术上的交流与指导，还是日常相处中的鼓励与陪伴，都让我在这段求学时光中感受到集体的温暖。那些一起聚餐、为彼此庆祝生日的平凡时光，早已悄然融入这段岁月，成为心中最温润的一隅。

寸草春晖，铭感于心。我还要特别感谢我的家人，他们始终如一的理解、支持与牵挂，是我求学路上最坚实的后盾。父亲默默的陪伴与坚定的支持、母亲细致入微的叮嘱与关心，让我在离家求学的岁月中始终感受到温暖与力量。正是有了家人的理解与付出，我才能更加从容地走过这段求学之路。

道阻且长，行则将至。同时，我也想感谢始终没有放弃的自己。或许我并不是天资出众的人，但在这段求学经历中，我始终以认真负责的态度对待学习、科研与生活，尽己所能，精益求精。一路走来，虽然有过迷茫与疲惫，但也正是在不断坚持与调整中，我逐渐学会了面对压力、承担责任，并在一次次尝试中获得成长。

行文至此，感念常在。多年以后，当我再度回望这段时光，依然会想起这段岁月里发生的点点滴滴，想起那些与我一同经历、彼此陪伴的人。三年的研究生生活虽将结束，但共同走过的日子不会因毕业而轻易消散，那些关于成长、陪伴与离别的记忆，终将沉淀于往后的岁月，并在某个再次提及的瞬间，重新变得清晰而温暖。千言万语，难尽此刻的不舍与感激。我会珍藏这些回忆，带着它们继续前行，也期待在人生更广阔的道路上，与大家以更好的模样再次相逢。谨以此文，向所有给予我关心、帮助与陪伴的师长、同学、朋友和家人，致以最诚挚的谢意。

